\documentclass{article}

% if you need to pass options to natbib, use, e.g.:
%     \PassOptionsToPackage{numbers, compress}{natbib}
% before loading neurips_2025

% Use this for double-blind submission:
\usepackage[dblblindworkshop]{neurips_2025}

% Use this for the camera-ready version:
% \usepackage[dblblindworkshop,final]{neurips_2025}

\workshoptitle{MATH-AI}

% Use this for preprint:
% \usepackage[preprint]{neurips_2025}

% to avoid loading the natbib package, add option nonatbib:
%    \usepackage[nonatbib,dblblindworkshop]{neurips_2025}

\usepackage[utf8]{inputenc} % allow utf-8 input
\usepackage[T1]{fontenc}    % use 8-bit T1 fonts
\usepackage{hyperref}       % hyperlinks
\usepackage{url}            % simple URL typesetting
\usepackage{booktabs}       % professional-quality tables
\usepackage{amsfonts}       % blackboard math symbols
\usepackage{nicefrac}       % compact symbols for 1/2, etc.
\usepackage{microtype}      % microtypography
\usepackage{xcolor}         % colors
\usepackage{graphicx}       % for including figures
\usepackage{amsmath}        % for math environments
\usepackage{amssymb}        % for additional math symbols

% Note. For the workshop paper template, both \title{} and \workshoptitle{} are required, with the former indicating the paper title shown in the title and the latter indicating the workshop title displayed in the footnote. 
\title{Your Paper Title Here}

% The \author macro works with any number of authors. There are two commands
% used to separate the names and addresses of multiple authors: \And and \AND.
%
% Using \And between authors leaves it to LaTeX to determine where to break the
% lines. Using \AND forces a line break at that point. So, if LaTeX puts 3 of 4
% authors names on the first line, and the last on the second line, try using
% \AND instead of \And before the third author name.

\author{%
  Author Name\thanks{Use footnote for providing further information
    about author (webpage, alternative address)---\emph{not} for acknowledging
    funding agencies.} \\
  Department/Institution \\
  University/Organization \\
  City, State/Country \\
  \texttt{email@domain.com} \\
  % Add more authors as needed:
  % \And
  % Second Author \\
  % Affiliation \\
  % Address \\
  % \texttt{email2@domain.com} \\
  % \AND
  % Third Author \\
  % Affiliation \\
  % Address \\
  % \texttt{email3@domain.com} \\
}

\begin{document}

\maketitle

\begin{abstract}
  Your abstract should be a concise summary of your work, typically 150-250 words. 
  It should clearly state the problem, your approach, key results, and implications. 
  The abstract must be limited to one paragraph and should be indented \nicefrac{1}{2}~inch 
  (3~picas) on both the left- and right-hand margins. Use 10~point type, with a vertical 
  spacing (leading) of 11~points.
\end{abstract}

\section{Introduction}

Start your introduction here. This section should:

\begin{itemize}
\item Motivate the problem and its importance
\item Provide necessary background and context
\item Clearly state your contributions
\item Outline the structure of the paper
\end{itemize}

Since this is a MATH-AI workshop submission, focus on the mathematical aspects and AI applications of your work.

\section{Related Work}

Discuss relevant prior work in this section. Since this is a double-blind submission, refer to your own work in the third person (e.g., "In previous work by Smith et al. [1]..." rather than "In our previous work [1]...").

\section{Methodology}

\subsection{Problem Formulation}

Clearly define the mathematical problem you are addressing. Use proper mathematical notation:

\begin{equation}
\min_{x \in \mathcal{X}} f(x) \quad \text{subject to} \quad g_i(x) \leq 0, \quad i = 1, \ldots, m
\end{equation}

\subsection{Your Approach}

Describe your proposed method in detail. Include:

\begin{itemize}
\item Mathematical foundations
\item Algorithmic details
\item Key theoretical results
\end{itemize}

\section{Experiments}

\subsection{Experimental Setup}

Describe your experimental setup, including:
\begin{itemize}
\item Datasets used
\item Evaluation metrics
\item Implementation details
\item Hyperparameter settings (can be moved to appendix)
\end{itemize}

\subsection{Results}

Present your main results. Use tables and figures effectively:

\begin{table}[h]
  \caption{Main experimental results}
  \label{tab:results}
  \centering
  \begin{tabular}{lcc}
    \toprule
    Method & Accuracy & Time (s) \\
    \midrule
    Baseline & 0.85 & 10.2 \\
    Your Method & 0.92 & 12.5 \\
    \bottomrule
  \end{tabular}
\end{table}

\begin{figure}[h]
  \centering
  \fbox{\rule[-.5cm]{0cm}{4cm} \rule[-.5cm]{4cm}{0cm}}
  \caption{Sample figure caption. Replace with your actual figure.}
  \label{fig:sample}
\end{figure}

\section{Discussion}

Discuss the implications of your results, limitations, and potential future work.

\section{Conclusion}

Summarize your contributions and key findings. Keep this section concise.

\begin{ack}
Use unnumbered first level headings for the acknowledgments. All acknowledgments
go at the end of the paper before the list of references. Moreover, you are required to declare
funding (financial activities supporting the submitted work) and competing interests (related financial activities outside the submitted work).
More information about this disclosure can be found at: \url{https://neurips.cc/Conferences/2025/PaperInformation/FundingDisclosure}.

Do {\bf not} include this section in the anonymized submission, only in the final paper. You can use the \texttt{ack} environment provided in the style file to automatically hide this section in the anonymized submission.
\end{ack}

\section*{References}

References follow the acknowledgments in the camera-ready paper. Use unnumbered first-level heading for
the references. Any choice of citation style is acceptable as long as you are
consistent. It is permissible to reduce the font size to \verb+small+ (9 point)
when listing the references.
Note that the Reference section does not count towards the page limit.
\medskip

{
\small

[1] Smith, J.\ \& Doe, A.\ (2023) Example paper title. In {\it Proceedings of Example Conference}, pp.\ 123--456.

[2] Johnson, B.\ (2022) Another relevant work. {\it Journal of Example Research} {\bf 15}(3):789-801.

[3] Wilson, C., Brown, D.\ \& Davis, E.\ (2021) Mathematical foundations for AI. {\it Example Journal} {\bf 42}(7):1234-1256.

}

%%%%%%%%%%%%%%%%%%%%%%%%%%%%%%%%%%%%%%%%%%%%%%%%%%%%%%%%%%%%

\appendix

\section{Technical Appendices and Supplementary Material}
Technical appendices with additional results, figures, graphs and proofs may be submitted with the paper submission before the full submission deadline (see above), or as a separate PDF in the ZIP file below before the supplementary material deadline. There is no page limit for the technical appendices.

\subsection{Additional Experimental Results}

Include additional experimental results, hyperparameter sensitivity analysis, or other supplementary material here.

\subsection{Proofs}

Include detailed proofs of theoretical results here.

%%%%%%%%%%%%%%%%%%%%%%%%%%%%%%%%%%%%%%%%%%%%%%%%%%%%%%%%%%%%

\end{document}
